Auditory attention decoding could let earbuds know which voice their wearer follows. But an earbud cannot be calibrated for each person, its two sides share no electrical reference, it must be causal, and it must work for people and speech it never met. Under those four conditions throughout, on public recordings, a learned decoder over four electrodes around each ear names the followed story in \nnSixty{} of one minute decisions and \nnFiveMin{} of five minute ones for 36 people it never saw (a linear decoder: \linSixty{}; chance: 50\%), keeps its accuracy on story content and on a source study it never trained on, and falls to chance when the speech is moved out of step. On 15 other people recorded in the ear, around it and on the scalp at the same moments, dry contacts in the ear canal decode at chance (\earInNn{}) while contacts around the ear reach \earAroundNn{}: an earbud's contacts must reach outside the canal. Replaying the real decisions under a head that looks away shows where the EEG adds to head direction, which is when the listener looks at someone other than the person followed and when attention moves while the head stays, and that the head should be weighed less than a hand set default. Degrading the raw recordings sets the electronics' tolerances. A jaw clench detector, run unchanged on four public datasets, catches \clAround{} of cued clenches around the ear, and a double clench fires by itself \clListenDouble{} times an hour in quiet listening against \clListenOthers{} for singles and holds, which makes the double clench the command that must not fire by accident.
